% Sample for pcalendar: every line shows the call (in monospace) followed by what it prints.
% Build: make test   (XeLaTeX; xepersian and the font Parisa are needed only to print the Persian text)
\documentclass[12pt]{article}

\usepackage[a4paper,margin=2cm]{geometry}
\usepackage{pcalendar}
\usepackage{pgffor}
\usepackage{xepersian}

\settextfont{Parisa}
\setlatintextfont{Latin Modern Roman}

\newcount\jd \newcount\wd

% \note{English text}: explanatory paragraph
\newcommand\note[1]{\par\medskip\noindent\lr{#1}\par\smallskip}
% \code{...}: the call, shown literally
\newcommand\code[1]{\lr{\texttt{\detokenize{#1}}}}
% \demo{call as text}{result}{comment}: one line = call, result, comment
\newcommand\demo[3]{\par\noindent\code{#1}\quad$\Rightarrow$\quad#2\ifx\relax#3\relax\else\quad\lr{(#3)}\fi}

% The date macros take a literal Y-M-D (a macro such as \d would not be expanded),
% so the dates are spelled out through these helpers rather than a \foreach list.
\newcommand\GtoJ[2]{\demo{\pdate{#1}}{\pdate{#1}}{}\demo{\pisodate{#1}}{\lr{\pisodate{#1}}}{#2}\par\smallskip}
\newcommand\JtoG[2]{\demo{\gdate{#1}}{\gdate{#1}}{}\demo{\gisodate{#1}}{\lr{\gisodate{#1}}}{#2}\par\smallskip}
\newcommand\YEAR[1]{\demo{\gisodate{#1-01-01}}{\lr{\gisodate{#1-01-01}}}{start of Jalali year #1}}
% one table cell: year and its animal; \YROW gives a row of years 6 apart
\newcommand\yc[1]{\edef\yy{\number\numexpr#1\relax}\lr{\yy}\quad\pcalendaryearname{\yy}}
\newcommand\YROW[1]{\yc{#1}&\yc{#1+6}&\yc{#1+12}&\yc{#1+18}}
\newcommand\WEEKDAY[2]{%
  \pcalendardatetojulian{#1}\jd
  \pgfcalendarjuliantoweekday{\jd}{\wd}%
  \advance\wd by 2 \ifnum\wd>6 \advance\wd by -7 \fi
  \par\noindent\code{#1}\quad$\Rightarrow$\quad\pcalendarweekdayname{\wd}\quad\lr{(#2)}}

\begin{document}
\setLTR % the explanations are English; only the results are Persian

\section*{\lr{pcalendar: sample and manual test}}

\note{pcalendar converts dates between the Persian (Jalali) and Gregorian calendars and provides
Persian month, weekday and year names. Each line below shows the call, an arrow, and the typeset result.
Only printing the Persian text needs a Unicode engine and a Persian font; this sample uses XeLaTeX with
\texttt{xepersian}. The conversions themselves work with any engine.}

\subsection*{\lr{1. Today}}
\note{\texttt{\textbackslash year}, \texttt{\textbackslash month} and \texttt{\textbackslash day} are
LaTeX's current date (Gregorian), so these lines change every day.}
\demo{\pdate{\year-\month-\day}}{\pdate{\year-\month-\day}}{today as day, month name, year}
\demo{\pisodate{\year-\month-\day}}{\lr{\pisodate{\year-\month-\day}}}{today as DD/MM/YYYY}

\subsection*{\lr{2. Gregorian to Jalali}}
\note{\texttt{\textbackslash pdate\{Y-M-D\}} takes a Gregorian date and prints the Jalali date with the month name.
\texttt{\textbackslash pisodate\{Y-M-D\}} prints it as DD/MM/YYYY, zero padded.}
\GtoJ{2015-03-21}{Nowruz 1394, the first day of the Jalali year}
\GtoJ{2021-03-20}{last day of a leap year, Esfand 30, 1399}
\GtoJ{2024-03-20}{Nowruz 1403}
\GtoJ{2024-12-31}{a day in Dey, which crosses the Gregorian new year}
\GtoJ{2025-03-20}{Esfand 30, 1403, the next leap-year end}

\subsection*{\lr{3. Jalali to Gregorian}}
\note{\texttt{\textbackslash gdate\{Y-M-D\}} and \texttt{\textbackslash gisodate\{Y-M-D\}} are the reverse:
the argument is a Jalali date, the result is Gregorian.}
\JtoG{1394-01-01}{Nowruz 1394}
\JtoG{1399-12-30}{last day of leap year 1399}
\JtoG{1403-01-01}{Nowruz 1403}
\JtoG{1403-10-11}{a day in Dey 1403}
\JtoG{1403-12-30}{last day of leap year 1403}

\subsection*{\lr{4. Start of each Jalali year}}
\note{The Gregorian date on which every fifth Jalali year starts. It is March 20, 21 or 22.}
\YEAR{1380}
\YEAR{1385}
\YEAR{1390}
\YEAR{1395}
\YEAR{1400}
\YEAR{1405}

\subsection*{\lr{5. Name macros}}
\note{These take a number and print the name. Months are 1 to 12; weekdays are 0 to 6,
where 0 is Saturday, the first day of the Persian week.}
\demo{\pcalendarmonthname{1}}{\pcalendarmonthname{1}}{Jalali month 1}
\demo{\pcalendarmonthname{12}}{\pcalendarmonthname{12}}{Jalali month 12}
\demo{\pcalendargregorianmonthname{3}}{\pcalendargregorianmonthname{3}}{Gregorian month 3, in Persian}
\demo{\pcalendarweekdayname{0}}{\pcalendarweekdayname{0}}{weekday 0}
\demo{\pcalendarweekdayname{6}}{\pcalendarweekdayname{6}}{weekday 6}
\demo{\pcalendarweekdayshortname{0}}{\pcalendarweekdayshortname{0}}{short weekday 0}
\par\medskip\noindent\lr{All Jalali months:}\quad
\foreach \m in {1,...,12}{\pcalendarmonthname{\m}\space}

\subsection*{\lr{6. Julian day numbers and date arithmetic}}
\note{\texttt{\textbackslash pcalendardatetojulian\{Y-M-D\}\{\textbackslash count\}} stores the Julian day
number of a Jalali date in a count register. A \texttt{+N} suffix adds N days.
\texttt{\textbackslash pcalendarjuliantodate\{\textbackslash count\}\{\textbackslash y\}\{\textbackslash m\}\{\textbackslash d\}}
converts a day number back and sets \texttt{\textbackslash y}, \texttt{\textbackslash m}, \texttt{\textbackslash d}.}
\pcalendardatetojulian{1394-01-01}\jd%
\demo{\pcalendardatetojulian{1394-01-01}\jd}{\number\jd}{day number of 1394-01-01}
\pcalendardatetojulian{1394-01-01+10}\jd%
\demo{\pcalendardatetojulian{1394-01-01+10}\jd}{\number\jd}{ten days later}
\pcalendarjuliantodate\jd\y\m\d%
\demo{\pcalendarjuliantodate\jd\y\m\d}{\lr{\number\y-\number\m-\number\d}}{back to a Jalali date}
\pcalendardatetojulian{1403-12-29+1}\jd%
\pcalendarjuliantodate\jd\y\m\d%
\demo{\pcalendardatetojulian{1403-12-29+1}\jd}{\lr{\number\y-\number\m-\number\d}}{1403 is a leap year, so Esfand has 30 days}
\pcalendardatetojulian{1403-12-29+2}\jd%
\pcalendarjuliantodate\jd\y\m\d%
\demo{\pcalendardatetojulian{1403-12-29+2}\jd}{\lr{\number\y-\number\m-\number\d}}{rolls over into the new year}

\subsection*{\lr{7. Weekday of a Jalali date}}
\note{The package has no weekday macro of its own. PGF's \texttt{\textbackslash pgfcalendarjuliantoweekday}
numbers Monday as 0, so the sample shifts the result by 2 to make Saturday 0, then calls
\texttt{\textbackslash pcalendarweekdayname}.}
\WEEKDAY{1394-01-01}{Saturday, March 21, 2015}
\WEEKDAY{1403-01-01}{Wednesday, March 20, 2024}
\WEEKDAY{1404-01-01}{Friday, March 21, 2025}

\subsection*{\lr{8. Animal year names}}
\note{\texttt{\textbackslash pcalendaryearname\{year\}} gives the animal of the 12-year cycle of the
Turkish-Mongol calendar for a Jalali year. It is defined for 0 to 2099 only.}
\medskip
\begin{tabular}{llll}
\YROW{1380}\\
\YROW{1381}\\
\YROW{1382}\\
\YROW{1383}\\
\YROW{1384}\\
\YROW{1385}\\
\end{tabular}

\subsection*{\lr{Limits}}
\note{The date arguments must be literal \texttt{Y-M-D} text. A macro holding a date, such as
\texttt{\textbackslash pdate\{\textbackslash mydate\}}, is not expanded and fails. The Jalali leap rule is the
33-year arithmetic approximation, which can differ from the official calendar in some years.}

\end{document}
